% =============================================================================
%  Robótica (MEEC, IST) — Lab 1: simulação do swing de golfe em MuJoCo
%  Compilar com:  latexmk -pdf main.tex   (corre pdflatex + bibtex as vezes
%  necessárias; no VS Code basta gravar, o LaTeX Workshop faz o resto)
%  Figuras em Photos/, geradas pelos scripts em ../../src/ (ver README.md):
%    v17_sequencia_face_on.png, v17_sequencia_down_the_line.png  <- renders.py --strip
%    v17_sequencia.png                                            <- sequencia.py
%    v17_dispersao.png                                            <- perturbacoes.py
%    v17_aerodinamica.png                                         <- teste_aerodinamica.py
%  Bibliografia: refs.bib (bibtex, estilo IEEEtran)
% =============================================================================
\documentclass[11pt,a4paper]{article}

\usepackage[utf8]{inputenc}
\usepackage[T1]{fontenc}
\usepackage[portuguese]{babel}
\usepackage[margin=2.5cm]{geometry}
\usepackage{graphicx}
\usepackage{booktabs}
\usepackage{tabularx}
\usepackage{amsmath}
\usepackage{siunitx}
\usepackage{caption}
\usepackage{xcolor}
\usepackage{listings}
\usepackage[hidelinks]{hyperref}

\sisetup{output-decimal-marker={,}, per-mode=symbol}

\lstset{basicstyle=\ttfamily\footnotesize, breaklines=true,
        frame=single, framesep=4pt, backgroundcolor=\color{gray!8},
        columns=flexible, keepspaces=true}

% \fig{ficheiro}{largura}: inclui a figura se existir; senão deixa um aviso
% visível, para que a compilação nunca falhe por uma figura ainda não gerada.
\newcommand{\fig}[2]{\IfFileExists{#1}{\includegraphics[width=#2]{#1}}%
  {\fbox{\parbox{#2}{\centering\ttfamily\small figura em falta: \detokenize{#1}}}}}

\title{\textbf{Golf swing in MuJoCo}\\[2mm]
       \large Robótica --- Laboratório 1}
\author{{Afonso Cruz (102815)} \and {Eduardo Batista (99931)} \and {Henrique Póvoa (102109)}}
\date{Grupo {nº\textbf{1}} \quad --- \quad 9/10/2026}

\begin{document}
\maketitle

\begin{abstract}
\noindent
Simula-se o swing de um jogador de golfe com um ferro~7, tratando o conjunto
jogador\,+\,taco como um robot articulado em MuJoCo~\cite{todorov2012}. O
modelo entregue (V17) tem três graus de liberdade atuados --- rotação do tronco,
elevação do braço esquerdo e pulso --- e um braço direito passivo preso ao
punho do taco. É controlado por binário calculado sobre uma trajetória de
polinómios de 5.º grau em que as três juntas cruzam a pose de endereço no
mesmo instante. O swing atinge \SI{28.2}{\meter\per\second} na cabeça do taco
e \SI{33.9}{\meter\per\second} na bola (\emph{smash factor} 1,20), com
binários dentro de limites humanos e a sequência proximal$\to$distal de um
swing real; a bola voa com arrasto e efeito de Magnus validados contra
coeficientes publicados. O movimento é anatomicamente plausível (penetração
máxima \SI{4.9}{\milli\meter}), mas fica abaixo da velocidade de um amador
(30--38\,\si{\meter\per\second}) e o ângulo de lançamento sai alto. Ensaios
com ruído de binário em cada junta mostram que o pulso é a junta sensível. Uma
versão anterior com dois graus de liberdade (V13), mais rápida mas com os
braços a atravessar o peito, serve de comparação. Usou-se MuJoCo 3.13 com
Python~3.11; todos os números deste relatório são reproduzidos pelos scripts
em \texttt{src/}.
\end{abstract}


% =============================================================================
\section{Tarefa 1 --- Modelo do jogador e do taco}
% =============================================================================

O modelo (\texttt{models/golfer\_v17.xml}) representa um adulto de
\SI{1.78}{\meter} e \SI{82}{\kilo\gram}, destro, na postura de endereço com um
ferro~7. Pernas e pélvis estão fixas ao mundo; o resto do corpo move-se com os
graus de liberdade da Tabela~\ref{tab:dof}.

\begin{table}[h]
\centering
\caption{Graus de liberdade do modelo.}
\label{tab:dof}
\small
\begin{tabularx}{\textwidth}{@{}l>{\raggedright\arraybackslash}Xcc@{}}
\toprule
Junta & Movimento & Amplitude & Binário máx. \\
\midrule
\multicolumn{4}{@{}l}{\emph{Atuadas}} \\
\texttt{torso}    & rotação do tronco em torno da coluna (inclinada \ang{35}) & \ang{-100} a \ang{95} & \SI{150}{\newton\meter} \\
\texttt{shoulder} & elevação do braço esquerdo à frente do peito, em torno da linha dos ombros & \ang{-115} a \ang{10} & \SI{100}{\newton\meter} \\
\texttt{wrist}    & taco em relação ao braço esquerdo, no plano do swing & \ang{-110} a \ang{120} & \SI{40}{\newton\meter} \\
\midrule
\multicolumn{4}{@{}l}{\emph{Passivas}} \\
\texttt{shoulder\_r}, \texttt{elbow\_r} & braço direito (rótula + charneira), mão presa ao punho & --- & --- \\
\texttt{neck}     & cabeça, acoplada a $-$\texttt{torso} & --- & --- \\
\texttt{ball\_free} & bola (6 dof) & --- & --- \\
\bottomrule
\end{tabularx}
\end{table}

\subsection{Juntas atuadas}

\begin{description}
  \item[\texttt{torso}] roda o tronco em torno da coluna. Os dois ombros são
        rígidos com o tronco, por isso o tronco leva os braços consigo. É a
        fonte de energia do swing: num corpo real ela vem das pernas e das
        ancas, aqui está toda neste motor (daí os \SI{150}{\newton\meter}: um
        ombro isolado produz 60--85\,\si{\newton\meter}~\cite{harbo2012}).
  \item[\texttt{shoulder}] levanta o braço esquerdo, que está sempre esticado.
        O eixo é a linha dos ombros: o braço sobe pela frente do peito, e é o
        tronco que o leva para trás, como num corpo humano. Esta escolha do
        eixo tem uma consequência geométrica importante para o braço direito
        (Secção~\ref{sec:passivas}).
  \item[\texttt{wrist}] roda o taco em relação ao braço esquerdo, no plano do
        swing. O limite de \SI{40}{\newton\meter} está acima da flexão máxima
        do pulso (22--25\,\si{\newton\meter}~\cite{harbo2012}) mas na ordem do
        binário que as mãos aplicam de facto ao taco num swing medido
        ($\approx$\SI{31}{\newton\meter}~\cite{nesbit2005}); a libertação é em
        boa parte passiva, centrífuga~\cite{sprigings2000}.
\end{description}

O braço esquerdo e o taco formam um pêndulo duplo, o modelo clássico do swing
de Cochran \& Stobbs e Jorgensen~\cite{cochran1968,jorgensen1970}. O pulso arma
no topo do backswing e só abre perto do impacto; atrasar essa libertação
aumenta a velocidade da cabeça~\cite{pickering1999}.

\subsection{Juntas passivas}
\label{sec:passivas}

O braço direito não tem motores. A mão direita está presa ao punho do taco por
uma restrição \texttt{connect}, e o braço limita-se a acompanhar o movimento:
o cotovelo dobra no backswing e estica no impacto. A restrição deixa um grau de
liberdade livre (o cotovelo pode rodar em torno da reta ombro--mão), que é
fixado por uma mola fraca na rótula (\SI{3}{\newton\meter\per\radian}).

Como a mão direita está presa ao punho, o braço direito tem de \emph{alcançar}
o punho em todas as poses. Se o eixo do ombro esquerdo passa pelos dois ombros,
a distância da mão esquerda ao ombro direito não muda com \texttt{shoulder} nem
com \texttt{torso} (ambos rodam os dois pontos em bloco), e o braço direito só
tem de absorver o movimento do pulso: a distância ombro-direito--punho é
\SI{0.60}{\meter} no endereço e \SI{0.54}{\meter} no topo, sempre abaixo do
braço esticado (\SI{0.62}{\meter}). Com outro eixo a restrição esticava no topo
e, por ser mole, puxava o taco. No swing nominal o erro máximo da restrição é
de \SI{2.2}{\milli\meter}.

A cabeça está ligada ao tronco por uma restrição entre juntas que a faz rodar
o mesmo ângulo em sentido contrário (\texttt{neck} $= -$\texttt{torso}). Assim
fica parada, a olhar para a bola. Nenhum destes corpos entra no modelo
dinâmico do controlador (Secção~\ref{sec:controlo}).

\subsection{Referenciais e convenção de sinais}

O referencial do mundo tem origem na bola, com $x$ na direção do alvo e $z$
para cima. O jogador está do lado $+y$. Nenhum corpo tem rotação própria em
relação ao pai (só translação), pelo que na pose de endereço todas as juntas
estão em $q=0$ e os eixos de todos os corpos coincidem com os do mundo: as
coordenadas dos \texttt{geom} lêem-se diretamente como coordenadas do mundo.
Os referenciais de cada corpo são: \texttt{torso} na anca; \texttt{head} no
topo do peito (um ponto sobre o eixo da coluna, para que a cabeça fique
exatamente parada); \texttt{arm\_l} no ombro esquerdo; \texttt{club} no punho
(mão esquerda); \texttt{upperarm\_r} e \texttt{forearm\_r} no ombro e no
cotovelo direitos.

A convenção de sinais é:
\begin{itemize}
  \item \textbf{tronco e pulso:} $q<0$ no backswing, $q>0$ no follow-through;
  \item \textbf{ombro:} $q<0$ com o braço levantado, tanto no backswing como no
        follow-through.
\end{itemize}

Daqui resulta a condição que domina todo o controlo: a cabeça do taco só passa
na bola quando as três juntas estão em $q=0$ ao mesmo tempo.

\subsection{Parâmetros físicos}

As massas seguem tabelas antropométricas aproximadas. As inércias são calculadas
pelo MuJoCo a partir da forma e da massa de cada \texttt{geom}.

\begin{itemize}
  \item \textbf{Braços:} \SI{4}{\kilo\gram} cada (2,2 + 1,3 + 0,5).
  \item \textbf{Taco:} \SI{0.43}{\kilo\gram}, dos quais \SI{0.26}{\kilo\gram}
        na cabeça; \emph{loft} de \ang{32}; haste rígida~\cite{milne1992}.
  \item \textbf{Bola:} \SI{45.9}{\gram}, \SI{42.7}{\milli\meter} de diâmetro,
        \SI{6}{\milli\meter} à frente da face no endereço (a mão direita presa
        por uma restrição mole faz o taco assentar uns milímetros; sem folga
        encostava na bola antes de o swing começar).
  \item \textbf{Binários máximos:} Tabela~\ref{tab:dof}; ver discussão acima.
  \item \textbf{Atrito nas juntas:} amortecimento viscoso de 2, 0{,}5 e
        \SI{0.1}{\newton\meter\second} no tronco, ombro e pulso; 0{,}5 nas
        juntas passivas do braço direito.
  \item \textbf{Inércia de rotor:} \SI{0.01}{\kilo\gram\meter\squared} em cada
        junta, para estabilizar a integração.
  \item \textbf{Ar:} densidade \SI{1.204}{\kilo\gram\per\meter\cubed} e
        viscosidade \SI{1.81e-5}{\pascal\second}, aplicado a todos os corpos
        (Secção~\ref{sec:aero}).
\end{itemize}

\subsection{Contacto e passo de integração}
\label{sec:contacto}

Só a cabeça do taco e a bola colidem. O corpo do jogador é apenas visual, e o
taco não toca no chão, o que evita ter de modelar o \emph{divot}. O impacto é
definido num par de contacto dedicado:

\begin{lstlisting}[language=XML]
<pair geom1="clubhead" geom2="ball_geom" condim="3"
      solref="0.0004 0.15" solimp="0.95 0.99 0.001"
      friction="0.4 0.4 0.005 0.0001 0.0001"/>
\end{lstlisting}

Os dois valores de \texttt{solref} têm significado físico:
\begin{itemize}
  \item \SI{0.4}{\milli\second} é a constante de tempo do contacto, igual à
        duração real de um impacto de golfe~\cite{penner2003};
  \item 0,15 é a razão de amortecimento. Sendo muito inferior a 1, o contacto
        fica sub-amortecido e a bola ressalta na face. É assim que se impõe a
        restituição no modelo de contacto mole do MuJoCo, que não tem
        coeficiente de restituição explícito~\cite{todorov2012}; o valor foi
        calibrado para \emph{smash factor} $\approx 1{,}2$.
\end{itemize}

Estes parâmetros estão no par e não no \texttt{geom} da bola. Se estivessem na
bola, seriam também usados no contacto com a relva.

O passo de integração é \SI{5e-5}{\second} no swing e no voo, mas desce para
\SI{5e-6}{\second} enquanto a cabeça do taco está a menos de \SI{15}{\centi\meter}
da bola. O contacto dura $\sim$\SI{0.4}{\milli\second}: com \SI{5e-5}{\second}
eram oito passos para o resolver, e a velocidade da bola variava $\pm5\,\%$
conforme o ponto do passo em que o contacto começava (uma mudança de
\SI{1}{\milli\second} na duração do downswing mudava o \emph{carry} em
dezenas de metros). Com o passo fino a variação cai para $\sim$1\,\%, e como
só dura alguns milissegundos de simulação o custo é desprezável. Pelo mesmo
motivo a velocidade da cabeça lê-se no estado \emph{anterior} ao passo em que o
contacto é detetado (depois dele já inclui parte do impulso), e a da bola
amostra-se também no passo em que o contacto acaba.

\subsection{Simplificações e impacto no realismo}

\begin{table}[h]
\centering
\caption{Simplificações e consequências.}
\label{tab:simpl}
\small
\begin{tabularx}{\textwidth}{@{}>{\raggedright\arraybackslash}p{4.6cm}X@{}}
\toprule
Simplificação & Impacto no realismo \\
\midrule
Pernas e pélvis rígidas &
Toda a rotação vem do tronco; não há rotação das ancas ($\sim\ang{45}$ num swing
real) nem transferência de peso~\cite{hume2005}. É a principal causa da
velocidade baixa da cabeça do taco. \\
Cotovelo esquerdo rígido &
O follow-through é curto: num swing real o braço da frente dobra depois do
impacto. \\
Ombro esquerdo com 1 dof &
Um ombro real tem 3 e a omoplata desliza. Com 1 dof o plano do swing é fixo. \\
Pulso com 1 dof, sem rotação dos antebraços &
A face do taco não abre nem fecha por ação do jogador, logo o modelo não
reproduz \emph{slice} nem \emph{hook}; a direção só varia por erro de
sincronização das juntas. \\
Braço direito passivo &
É só carga: não empurra o taco (num swing real o braço de trás contribui para
a velocidade). O seu efeito entra no controlador como perturbação. \\
Corpo sem contactos &
Nada impede fisicamente os braços de atravessar o corpo; a penetração máxima
foi medida à parte: \SI{4.9}{\milli\meter} (braço esquerdo no peito, perto do
topo). \\
Taco rígido, cabeça em caixa &
Sem flexão da haste (contributo pequeno~\cite{milne1992}) nem geometria real
da face: o ponto de impacto na face não é modelado. \\
Contacto mole, restituição calibrada &
O \emph{smash factor} sai do \texttt{solref}, não de um material. A bola sai
pela normal da face (Secção~\ref{sec:comentario}). \\
Aerodinâmica com coeficientes constantes &
Arrasto $C_d$ e sustentação $C_L$ não dependem do Reynolds nem saturam com a
rotação: o Magnus do MuJoCo é linear em $\omega$. Calibrado para a rotação
baixa deste swing, exagera o voo a 7000~rpm (Secção~\ref{sec:aero}). \\
Ar no corpo por elipsoides genéricos &
Cada \texttt{geom} do braço e do taco sente o ar como um elipsoide com os
coeficientes por omissão do MuJoCo. Efeito pequeno (\SI{1.8}{\newton\meter} no
pulso). \\
Relva como plano rígido &
O rolamento é exagerado (sem deformação nem travagem pelo backspin); usa-se o
\emph{carry} como métrica~\cite{penner2002}. \\
Motores de binário ideais &
Sem relação força--velocidade muscular nem atraso de ativação: o binário máximo
está disponível instantaneamente, também na fase mais rápida. \\
Controlador com a dinâmica exata &
O binário calculado usa $M(q)$ e $h(q,\dot q)$ do próprio modelo; um humano tem
um modelo interno aproximado e atrasos de $\sim$\SI{100}{\milli\second}. \\
\bottomrule
\end{tabularx}
\end{table}

% =============================================================================
\section{Tarefa 2 --- Controlo e ensaios}
% =============================================================================

\subsection{Trajetória de referência e controlador}
\label{sec:controlo}

A referência de cada junta é uma sequência de polinómios de 5.º grau
(backswing de \SI{0.80}{\second}, pausa de \SI{0.05}{\second}, downswing de
\SI{0.30}{\second} e follow-through), com aceleração nula nas extremidades de
cada troço, pelo que posição, velocidade e aceleração são contínuas e os
binários pedidos não têm degraus~\cite{sequeira2024}. O topo é
$(\ang{-90}, \ang{-85}, \ang{-95})$ para tronco, ombro e pulso. A trajetória
imita a sequência de um swing humano:

\begin{itemize}
  \item \textbf{Backswing:} tronco e braço arrancam juntos; o pulso só começa a
        armar a \SI{30}{\percent} do backswing.
  \item \textbf{Downswing:} o tronco arranca primeiro, o braço a
        \SI{10}{\percent}, e o pulso fica armado até \SI{55}{\percent}
        (\emph{lag}).
  \item \textbf{Impacto:} as três juntas chegam a $q=0$ no mesmo instante
        ($t=\SI{1.15}{\second}$), mas com velocidades diferentes: o braço chega
        parado (ponto mais baixo do arco das mãos), o tronco ainda a rodar
        (\SI{425}{\degree\per\second}) e o pulso no pico de velocidade. Para o
        pulso impõe-se $v_{\text{imp}} = 2|q_{\text{topo}}|/T$: com essa
        condição o polinómio do downswing reduz-se a
        $q = \Delta q\,(2\tau^3 - \tau^4)$, cuja velocidade é monótona e tem o
        máximo exatamente no impacto.
  \item \textbf{Follow-through:} cada junta trava até à pose final
        $(\ang{85}, \ang{-70}, \ang{105})$ num troço simétrico do downswing.
\end{itemize}

A trajetória é seguida por binário calculado (\emph{inverse dynamics
control}~\cite{siciliano2009,sequeira2024}),
\begin{equation}
  \boldsymbol\tau = M(q)\left[\ddot q_{\text{ref}}
      + K_d\left(\dot q_{\text{ref}} - \dot q\right)
      + K_p\left(q_{\text{ref}} - q\right)\right] + h(q,\dot q),
  \label{eq:ct}
\end{equation}
com $M$ a submatriz de massa das três juntas atuadas (\texttt{mj\_fullM}), $h$
os termos de Coriolis, centrífugos e de gravidade (\texttt{qfrc\_bias}),
$K_p = 900$ e $K_d = 60$. O braço direito e a cabeça não entram em $M$; o seu
efeito é tratado como perturbação e corrigido pelo PD. Os binários são saturados
nos limites do XML.

\subsection{Bateria de testes}

Para avaliar o realismo do modelo correram-se, por esta ordem: (i) um pêndulo
simples e um pêndulo a bater numa bola (dica~1 do enunciado; período com erro
$<0{,}2\,\%$ face à teoria); (ii) um estudo de convergência do passo de
integração no contacto (Secção~\ref{sec:contacto}); (iii) a bola em voo livre,
para medir os coeficientes aerodinâmicos que o MuJoCo aplica de facto
(Secção~\ref{sec:aero}); (iv) o swing nominal, comparado com médias publicadas
de um ferro~7 (Tabela~\ref{tab:nominal}); (v) a sequência cinemática das juntas
(Figura~\ref{fig:seq}); (vi) ruído de binário em cada grau de liberdade
(Secção~\ref{sec:pert}); (vii) o mesmo swing com três modelos de ar no corpo;
(viii) a ``tacada perfeita'': a bola lançada no mesmo modelo com as condições
publicadas de um ferro~7, para separar o que é do impacto do que é do voo.

\subsection{Ensaio nominal}
\label{sec:nominal}

\begin{figure}[h]
\centering
\fig{Photos/v17_sequencia_face_on.png}{\textwidth}\\[1mm]
\fig{Photos/v17_sequencia_down_the_line.png}{\textwidth}
\caption{Fases do swing simulado, câmaras \emph{face-on} (cima) e
\emph{down-the-line} (baixo).}
\label{fig:fases}
\end{figure}

A Figura~\ref{fig:fases} mostra as sete fases do swing. O braço direito dobra no
topo e estica no impacto sem ser comandado, a cabeça fica parada, e a pose de
impacto reproduz a de endereço. A penetração máxima medida entre partes móveis e
o corpo é de \SI{4.9}{\milli\meter} (braço esquerdo contra o peito, perto do
topo).

\begin{figure}[h]
\centering
\fig{Photos/v17_sequencia.png}{0.72\textwidth}
\caption{Ângulos, velocidades e binários das três juntas atuadas.}
\label{fig:seq}
\end{figure}

Na Figura~\ref{fig:seq} o seguimento da referência é quase exato, com as três
juntas a cruzar zero no impacto ($-1{,}0$, $+1{,}1$ e $+2{,}7$\,\si{\degree}).
As velocidades mostram o padrão proximal $\to$ distal: ombro e tronco atingem o
pico ($\num{598}$ e \SI{453}{\degree\per\second}) \SI{0.13}{} e
\SI{0.07}{\second} antes do impacto e travam, enquanto o pulso dispara para
\SI{1392}{\degree\per\second} exatamente no impacto. A ordem entre tronco e
braço não é a de um swing real (o braço atinge o pico antes do tronco): é uma
consequência de se impor ao braço velocidade nula no impacto. Os binários ficam
dentro dos limites: tronco e ombro nunca saturam no downswing (picos de 120 e
\SI{91}{\newton\meter}) e o pulso satura em \SI{13}{\percent} do downswing.

\begin{table}[h]
\centering
\caption{Swing nominal e médias publicadas de um ferro~7 (TrackMan, LPGA a PGA
Tour~\cite{trackman2024}; amador: \cite{golfmonthly_7iron}).}
\label{tab:nominal}
\small
\begin{tabular}{lccl}
\toprule
Grandeza & Modelo & Real & \\
\midrule
Duração do downswing          & \SI{0.30}{\second} & 0,26--0,39~\cite{grober2010} & ok \\
Juntas no impacto (tronco, ombro, pulso) & $(\ang{-1.0},\,\ang{+1.1},\,\ang{+2.7})$ & $(0,0,0)$ & ok \\
Velocidade da cabeça do taco  & \SI{28.2}{\meter\per\second} & 34,0--40,2 (amador $\sim$36) & baixo \\
Velocidade da bola            & \SI{33.9}{\meter\per\second} & 44,7--53,6 & baixo \\
\emph{Smash factor}           & 1,20 & 1,33--1,38 & baixo \\
Ângulo de ataque              & \ang{-1.6} & $\sim\ang{-4}$ & ok \\
Ângulo de lançamento          & \ang{33.6} & 16--19 & alto \\
\emph{Spin}                   & \SI{2510}{rpm} & 6700--7100 & baixo \\
\emph{Carry}                  & \SI{94.6}{\meter} & 129--157 & baixo \\
Altura máxima                 & \SI{21.1}{\meter} & 24--29 & ok \\
Ângulo de descida             & \ang{44.4} & 47--50 & ok \\
Desvio lateral na aterragem   & \SI{-0.7}{\meter} & --- & \\
\bottomrule
\end{tabular}
\end{table}

\subsection{Aerodinâmica}
\label{sec:aero}

O ar está ligado para todo o modelo (\texttt{density}, \texttt{viscosity}). A
bola usa o modelo de elipsoide do MuJoCo com arrasto e sustentação de
Magnus~\cite{mujoco_fluid}; como o MuJoCo não usa o fator $\tfrac12$ da
definição clássica, os valores de \texttt{fluidcoef} não são os coeficientes
da literatura e mediram-se em voo livre (Tabela~\ref{tab:aero}). O decaimento
da rotação foi ajustado (\texttt{fluidcoef} angular 0,1 em vez de 1,5, que
tirava \SI{34}{\percent} por segundo e anulava o Magnus a meio do voo).

\begin{table}[h]
\centering
\caption{Coeficientes aerodinâmicos efetivos da bola, medidos em voo livre nas
condições de lançamento do swing ($S = r\omega/v = 0{,}17$).}
\label{tab:aero}
\footnotesize
\begin{tabular}{lcc}
\toprule
 & Modelo & Literatura \\
\midrule
$C_d$ & 0,28 & 0,27--0,32~\cite{bearman1976,smits1994} \\
$C_L$ a $S\approx0{,}15$ & 0,22 & 0,15--0,25~\cite{bearman1976,smits1994} \\
Decaimento da rotação & \SI{2.8}{\percent\per\second} & 4--5\,\si{\percent\per\second} a \SI{45}{\meter\per\second}~\cite{smits1994,tavares1999} \\
\bottomrule
\end{tabular}
\end{table}

\begin{figure}[h]
\centering
\fig{Photos/v17_aerodinamica.png}{\textwidth}
\caption{Esquerda: voo da bola com as condições de saída do swing, no vácuo,
só com arrasto e com arrasto\,+\,Magnus. Direita: velocidade da cabeça do taco
com três modelos de ar no braço e no taco, com e sem compensação do arrasto
pelo controlador.}
\label{fig:aero}
\end{figure}

Dois resultados merecem registo. Primeiro, a validação cruzada: a bola lançada
sozinha com as condições de saída do swing aterra a \SI{94.6}{\meter}, o mesmo
valor do swing completo; sem Magnus daria \SI{78}{\meter} e no vácuo
\SI{108}{\meter} (Figura~\ref{fig:aero}). Segundo, o ar no corpo: o modelo que
o MuJoCo usa por omissão quando um \texttt{geom} não declara \texttt{fluidshape}
(uma ``caixa de inércia'' equivalente) dá \SI{11.8}{\newton\meter} de arrasto no
pulso, seis vezes mais do que o modelo por elipsoides (\SI{1.8}{\newton\meter})
e do que uma estimativa à mão com $C_d \approx 1{,}1$ para a cabeça e a
haste~\cite{hoerner1965} (\SI{1.5}{\newton\meter}): para um corpo comprido e
fino com a massa na ponta, uma caixa com a mesma inércia tem 3$\times$ a área
frontal. Por isso todos os \texttt{geom} declaram \texttt{fluidshape="ellipsoid"}.
Neste modelo o ar quase não altera a velocidade do taco (28,2 vs
28,2\,\si{\meter\per\second} no vácuo), porque tronco e ombro não saturam e o
PD compensa; no V13, com os motores saturados no fim do downswing, o ar tirava
1,3\,\%.

A limitação está na ``tacada perfeita'': lançada no mesmo modelo com as
condições médias do LPGA Tour (\SI{46.5}{\meter\per\second}, \ang{19},
\SI{6700}{rpm}), a bola voa \SI{162}{\meter} com \SI{44}{\meter} de altura,
contra 129 e \SI{24}{\meter} publicados. O Magnus do MuJoCo é linear na
rotação, enquanto numa bola real $C_L$ satura ($C_L \propto S^{0{,}4}$
em~\cite{smits1994}); o modelo está calibrado para a rotação baixa deste swing,
não para as 7000~rpm reais.

\subsection{Ensaios com perturbações}
\label{sec:pert}

Somou-se ruído gaussiano de $\sigma = \SI{10}{\newton\meter}$ ao binário de uma
junta de cada vez e repetiu-se o swing $N=20$ vezes por junta. Escolheu-se
ruído no binário, e não na posição, porque o ruído motor humano cresce com o
comando~\cite{harris1998}; é mantido constante durante \SI{0.1}{\second}
(\SI{10}{\hertz}, a banda do tremor fisiológico
reforçado~\cite{elble1990,deuschl1998}) em vez de reamostrado a cada passo: com
ruído branco por passo as $\sim$6000 amostras de um downswing cancelam-se e o
resultado passa a depender do passo de integração. O ruído só atua depois de o
taco sair do endereço ($t \geq \SI{0.24}{\second}$): com a face a milímetros
da bola, um desequilíbrio de \SI{10}{\newton\meter} no pulso em repouso dá um
erro estático do PD de $\tau/(K_p M_{ww}) \approx \ang{3}$, cerca de
\SI{35}{\milli\meter} na cabeça, e deslocava a bola antes de o swing começar.
Note-se que $\sigma$ é o mesmo em valor absoluto mas não em fração da
capacidade de cada motor (coluna $\sigma/\tau_{\max}$).

\begin{table}[h]
\centering
\caption{Efeito do ruído de binário ($\sigma=\SI{10}{\newton\meter}$, $N=20$),
média $\pm$ desvio padrão dos swings válidos. Lateral: $+$ = esquerda do alvo
(lado do jogador), $-$ = direita.}
\label{tab:pert}
\small
\begin{tabular}{lcccccc}
\toprule
Junta & $\sigma/\tau_{\max}$ & Válidos & $v_{\text{cab}}$ [\si{\meter\per\second}] & Lanç. [\si{\degree}] & \emph{Carry} [\si{\meter}] & Lateral [\si{\meter}] \\
\midrule
\texttt{torso}    & 7\,\%  & 20/20 & $28{,}2 \pm 0{,}0$ & $33{,}8 \pm 0{,}2$ & $95{,}3 \pm 0{,}4$ & $-0{,}8 \pm 0{,}2$ \\
\texttt{shoulder} & 10\,\% & 19/20 & $28{,}2 \pm 0{,}1$ & $33{,}7 \pm 0{,}1$ & $95{,}4 \pm 0{,}6$ & $-1{,}4 \pm 2{,}6$ \\
\texttt{wrist}    & 25\,\% & 20/20 & $28{,}1 \pm 0{,}4$ & $33{,}6 \pm 0{,}8$ & $94{,}7 \pm 2{,}3$ & $-0{,}5 \pm 0{,}8$ \\
\bottomrule
\end{tabular}
\end{table}

\begin{figure}[h]
\centering
\fig{Photos/v17_dispersao.png}{\textwidth}
\caption{Pontos de aterragem dos swings com ruído, um painel por junta
(estrela: nominal; barras: média $\pm$ 1 desvio padrão).}
\label{fig:dispersao}
\end{figure}

\textbf{O pulso é a junta mais sensível no alcance.} É a única em que o ruído
se vê na velocidade da cabeça ($\pm$\SI{0.4}{\meter\per\second}) e no
\emph{carry} ($\pm$\SI{2.3}{\meter}, de 90{,}7 a \SI{97.9}{\meter}): os
\SI{10}{\newton\meter} são \SI{25}{\percent} do motor do pulso e atuam sobre a
inércia pequena do taco, e os swings mais curtos são os em que o pulso chega
ao impacto ainda armado ($+4$ a \ang{+6} em vez de \ang{+2.7}). No tronco e no
ombro, com inércias e binários maiores, o mesmo ruído quase não altera o
alcance ($\pm$0,4 e $\pm$\SI{0.6}{\meter}).

\textbf{O ombro é a junta sensível na direção.} Dois swings com ruído no ombro
aterraram 5 e \SI{11}{\meter} para a direita do alvo e um falhou a bola. A
causa não está na orientação da face (a sincronização das juntas é igual à
dos outros swings) mas na altura das mãos: um erro de \ang{0.7} no ombro
desloca a cabeça do taco quase \SI{1}{\centi\meter} na vertical, e a bola é
apanhada fora do centro da face --- ou nem sequer apanhada. O modelo só
reproduz erros de direção por este mecanismo, porque sem rotação dos
antebraços a face não pode abrir nem fechar.

\textbf{Comparação com o V13.} Na versão de dois graus de liberdade o mesmo
ensaio dava, para o pulso, 4 em 20 bolas a 10--\SI{17}{\meter} do alvo e uma
``topada'', com o \emph{carry} a variar $\pm$\SI{32}{\meter}. O V17 rejeita o
ruído muito melhor por duas razões: o pulso só está saturado \SI{13}{\percent}
do downswing (contra \SI{30}{\percent} no V13), pelo que o PD tem autoridade
para corrigir, e o ruído só começa depois de o taco sair do endereço. A
primeira versão deste ensaio, com ruído desde $t=0$ e a face a
\SI{1.7}{\milli\meter} da bola, dava 11 falhas em 20 no pulso: o taco empurrava
a bola antes do swing. Esse resultado media o desenho do teste, não o modelo.

\subsection{Comparação com a versão de dois graus de liberdade}

\begin{table}[h]
\centering
\caption{V13 (braços em bloco + pulso) vs V17 (tronco + braço esquerdo + pulso,
braço direito passivo). Ambos com o mesmo taco, bola, contacto e ar.}
\label{tab:v13v17}
\small
\begin{tabular}{lcc}
\toprule
 & V13 (2 dof) & V17 (3 dof) \\
\midrule
Velocidade da cabeça do taco [\si{\meter\per\second}] & 38,8 & 28,2 \\
Bola [\si{\meter\per\second}] / \emph{smash} & 47,2 / 1,22 & 33,9 / 1,20 \\
Ângulo de lançamento [\si{\degree}] & 31,3 & 33,6 \\
\emph{Spin} [rpm] & 3200 & 2510 \\
\emph{Carry} [\si{\meter}] / altura máx. [\si{\meter}] & 151 / 39 & 95 / 21 \\
Juntas no impacto [\si{\degree}] & $(+1{,}1,\,-1{,}8)$ & $(-1{,}0,\,+1{,}1,\,+2{,}7)$ \\
Saturação no downswing & 20\,\% / 30\,\% & 0 / 0 / 13\,\% \\
Braços atravessam o peito & sim (vários cm) & não (\SI{4.9}{\milli\meter}) \\
Braço direito & rígido & passivo (dobra e estica) \\
\bottomrule
\end{tabular}
\end{table}

A Tabela~\ref{tab:v13v17} resume a troca feita: o V13 roda os dois braços em
bloco em torno do peito com \SI{150}{\newton\meter} e um downswing de
aceleração constante com os motores no limite, e chega a
\SI{38.8}{\meter\per\second} --- dentro da gama real --- mas com os braços a
atravessar o tronco. O V17 é anatomicamente plausível e perde
\SI{27}{\percent} de velocidade.

% =============================================================================
\section{Tarefa 3 --- Comentário aos resultados}
\label{sec:comentario}
% =============================================================================

\textbf{O que correu bem.} O movimento é o de um swing reconhecível
(Figura~\ref{fig:fases}): tronco a rodar, braço da frente esticado, pulso
armado no topo e libertado tarde, braço de trás a dobrar e esticar sozinho
(erro da restrição mão--punho $\leq$\SI{2.2}{\milli\meter}), cabeça parada. A
trajetória é seguida com binários dentro de limites humanos, sem degraus, e o
impacto ocorre a menos de \ang{3} da pose de endereço nas três juntas. A
sequência de picos de velocidade é a proximal$\to$distal dos swings
reais~\cite{hume2005}, com o pulso no máximo exatamente no impacto, como prevê
o pêndulo duplo~\cite{jorgensen1970,pickering1999}. O voo da bola usa
coeficientes aerodinâmicos dentro das gamas publicadas, verificados por
medição e por validação cruzada. E o resultado é numericamente repetível: o
estudo de convergência do contacto eliminou uma variação de $\pm5\,\%$ na
velocidade da bola que, numa versão anterior, se confundia com dispersão
devida ao ruído.

\textbf{O que não correu bem.} A cabeça do taco chega a
\SI{28.2}{\meter\per\second}, abaixo da gama de um amador. A causa não é falta
de binário --- tronco e ombro nunca saturam --- mas a trajetória imposta: o
braço chega parado ao impacto e só o tronco e o pulso contribuem para a
velocidade da cabeça, e com as ancas fixas falta a rotação e a transferência de
peso que num corpo real alimentam o tronco. O V13, com os dois braços rígidos a
rodar no peito, atingia \SI{38.8}{\meter\per\second}; trocou-se velocidade por
plausibilidade anatómica. O \emph{smash factor} de 1,20 também fica abaixo do
real (1,33), por ser um parâmetro de contacto calibrado e não uma propriedade
do material.

\textbf{Depois do impacto.} O ângulo de lançamento (\ang{33.6}) é praticamente
o \emph{loft} dinâmico da face (\ang{33.7}): a bola sai pela normal da face,
quando num impacto real o atrito tangencial a faz sair mais baixa (16--\ang{19})
e com muito mais rotação~\cite{penner2001,penner2003}. Há duas causas: o modelo
não tem o grau de liberdade que inclina a haste para a frente no impacto (as
mãos chegam na pose de endereço, com o \emph{loft} nominal), e o contacto mole
do MuJoCo, num impacto de \SI{0.4}{\milli\second}, não reproduz a bola a rolar
pela face. O \emph{carry} de \SI{95}{\meter} resulta da velocidade baixa; a
altura máxima e o ângulo de descida saem na gama real, mas por compensação
entre lançamento alto e rotação baixa, não porque o voo esteja certo --- a
tacada perfeita a 6700~rpm mostra que o Magnus linear do MuJoCo exagera a
sustentação a rotações reais.

\textbf{Perturbações.} O pulso é a junta sensível: controla a orientação da
face no impacto, e \SI{10}{\newton\meter} são \SI{25}{\percent} do seu motor.
O controlador conhece a dinâmica exata e tem ganhos altos, por isso rejeita o
ruído melhor do que um humano com atrasos de $\sim$\SI{100}{\milli\second}: a
dispersão medida é um limite inferior da real.

\textbf{Trabalho futuro.} Por ordem de efeito esperado: libertar as ancas (ou
dar velocidade ao braço no impacto) para ganhar velocidade; um grau de
liberdade que incline a haste no impacto, para baixar o \emph{loft} dinâmico;
a rotação dos antebraços, sem a qual o ensaio de perturbações não consegue
mostrar erros de direção por abertura da face; e um Magnus com saturação em
$S$, calculado em Python e aplicado como força externa.

% =============================================================================
\bibliographystyle{IEEEtran}
\bibliography{refs}

\end{document}
